% Vietnamese translation for OI Public Library
% Source-File: IOI中国国家候选队论文/2005/杨思雨/杨思雨.ppt
\documentclass[11pt]{article}
\usepackage{oipl-vn}

\title{Bản trình chiếu: Tỉ lệ vàng và tin học}
\author{Yang Siyu}
\date{2005}

\begin{document}
\maketitle

\origtitle{Slide companion for the golden ratio in informatics}
\sourcepath{IOI China candidate team papers / 2005 / Yang Siyu / Yang Siyu.ppt}

\section{Phạm vi bản dịch}

Bản trình chiếu đi cùng bài viết về liên hệ giữa tỉ lệ vàng và thi tin học.
Phần chữ đọc được có các mốc \(10^{-3}\), \(1/108\), \(2\log_2 108\), cùng các
hình nền nghệ thuật như tranh của Vincent van Gogh.

\section{Ý tưởng chính}

Tỉ lệ vàng xuất hiện trong các bài toán tối ưu một chiều, chia đoạn và trò
chơi lấy đá kiểu Wythoff. Khi lời giải phụ thuộc vào việc chia không gian tìm
kiếm sao cho số khả năng còn lại cân bằng, các hằng số gần tỉ lệ vàng có thể
tự nhiên xuất hiện.

\section{Trong thuật toán}

Tìm kiếm theo tỉ lệ vàng dùng hai điểm kiểm tra để giảm số lần đánh giá hàm
đơn đỉnh. Trong trò chơi tổ hợp, các cặp vị trí thua có thể liên hệ với
\(\varphi=(1+\sqrt5)/2\). Trong phân tích độ chính xác, các công thức logarit
trên slide cho thấy số bước cần thiết để đạt sai số cho trước.

\section{Kết luận}

Bản trình chiếu mang tính trực quan: nó dùng hình ảnh và ví dụ để nhấn mạnh
rằng cấu trúc toán học đẹp có thể dẫn đến thuật toán gọn, nhưng vẫn cần chứng
minh điều kiện áp dụng trong từng bài.

\end{document}
